% source: data/dce_emp_wtp.csv, data/dce_org_wtp.csv, data/dce_design.csv ; facts: E1-E7, R1-R7, G1-G6
% generated by deck_support/scripts/make_tables.py for B-gap; do not edit by hand
% Employer update (hand edit, appended sample n = 397): employer discounts from appended_wtp.xlsx sheet main.
% appended_wtp.xlsx has no gap sheet, so each gap = employer discount - PME value (EB E1-E7), with
% SE = sqrt(se_PME^2 + se_employer^2) (independent samples, as in employer_wtp.xlsx sheet gap), SEs taken
% from the 95% CIs; the employer values are rounded to whole points, so gaps are approximate (+/- 0.5).
% Change vs the 303-firm table: the 1-2 days gap is now statistically significant (13, 95% CI 2-23, p = .02;
% was 11, n.s., p = .12).
\begin{tabular}{@{}>{\raggedright\arraybackslash}p{3.30cm}@{\hspace{0.25cm}}>{\raggedleft\arraybackslash}p{1.65cm}@{\hspace{0.25cm}}>{\raggedright\arraybackslash}p{2.90cm}@{\hspace{0.25cm}}>{\raggedright\arraybackslash}p{2.85cm}@{\hspace{0.25cm}}>{\raggedright\arraybackslash}p{2.80cm}@{}}
\toprule
\multicolumn{1}{@{}>{\raggedright\arraybackslash}b{3.30cm}@{\hspace{0.25cm}}}{\textbf{Arrangement}} & \multicolumn{1}{@{}>{\raggedleft\arraybackslash}b{1.65cm}@{\hspace{0.25cm}}}{\textbf{PMEs' value}} & \multicolumn{1}{@{}>{\raggedright\arraybackslash}b{2.90cm}@{\hspace{0.25cm}}}{\textbf{Employers' model-implied discount}} & \multicolumn{1}{@{}>{\raggedright\arraybackslash}b{2.85cm}@{\hspace{0.25cm}}}{\textbf{Gap (95\% CI)}} & \multicolumn{1}{@{}>{\raggedright\arraybackslash}b{2.80cm}@{}}{\textbf{Zone}} \\
\midrule
Compressed week & \textcolor{Muted}{\ensuremath{\approx}7} & \textcolor{Muted}{\ensuremath{\approx}7} & \textcolor{Muted}{0 ($-$13 to 12), n.s.} & A: common ground \\
\addlinespace[3pt]
1--2 days WFH & \textcolor{Muted}{\ensuremath{\approx}7} & \textcolor{Muted}{\ensuremath{\approx}20} & \textcolor{Muted}{13 (2--23), p~=~.02} & B: depends on the role \\
\addlinespace[3pt]
Staggered hours & \textcolor{Muted}{\ensuremath{\approx}0} & \textcolor{Muted}{\ensuremath{\approx}10, \Pend{\emph{not robust}}} & \textcolor{Muted}{--} & Little at stake \\
\addlinespace[3pt]
3--4 days WFH & \textcolor{Muted}{\ensuremath{\approx}10} & \textcolor{Muted}{\ensuremath{\approx}36} & \textcolor{Muted}{27 (16--37)} & C: far apart \\
\addlinespace[3pt]
Fully remote & \textcolor{Muted}{\ensuremath{\approx}15\textsuperscript{1}} & \textcolor{Muted}{beyond the largest discount tested (33\%)} & \textcolor{Muted}{widest; \emph{not robust}} & C: far apart \\
\addlinespace[3pt]
Reduced load / job sharing & \textcolor{Muted}{\ensuremath{\approx}6} & \textcolor{Muted}{\ensuremath{\approx}24} & \textcolor{Muted}{\ensuremath{\approx}18} & Pay for hours \\
\addlinespace[3pt]
Part-time & \textcolor{Muted}{\ensuremath{\approx}4} & \textcolor{Muted}{\ensuremath{\approx}36} & \textcolor{Muted}{\ensuremath{\approx}32} & Pay for hours \\
\bottomrule
\addlinespace[3pt]
\multicolumn{5}{@{}>{\raggedright\arraybackslash}p{14.50cm}@{}}{\leavevmode\color{Muted}\fontsize{8.7}{9.8}\selectfont \textbf{Zones are descriptive; these numbers are indicative.} Gap = discount minus value. Pay gaps shown were 17--33\%, so zones rest on equal-pay choices. n.s. = not statistically significant at the 5\% level. \textsuperscript{1}vs onsite with a long commute. Grey = indicative.} \\
\end{tabular}
